Discussions of Generation Z and workplace ethics usually ask what this cohort values. This paper has asked a different question: how members of the cohort learned the ethics that are actually practiced where they work, given that many of them entered office-based work when distributed arrangements were spreading. Drawing on social learning theory and imprinting theory, it proposed that distributed entry lowers newcomers' clarity about enacted ethical norms by reducing their observational access to ethical conduct, that leaders who narrate the ethical reasoning behind their decisions can offset part of that loss, and that distance weakens the transmission of enacted norms whether those norms are ethical or not. It further proposed that the resulting differences outlast the entry period and that they follow the conditions of entry rather than birth cohort. The main takeaway is that the channel through which newcomers learn workplace ethics deserves as much attention as the values they bring with them. The main limitation is that the model is untested and three of its constructs still lack measures. A useful next step would be to compare distributed and co-located entrants within the same organizations, including older entrants, and to follow both groups after their working arrangements converge. Such a study would show whether the ethical imprint of distributed entry is real, whether it lasts, and whether it belongs to a generation or to a set of entry conditions that any newcomer might face.
